\section{Experiencia Profesional}
\vspace{-1.5em}
\textcolor{darkgray}{\rule{\textwidth}{0.5pt}}
\begin{longtable}{p{14cm}>{\raggedleft\arraybackslash}p{4cm}}
\textbf{Asistente de Investigación} & \textbf{Dic 2025 -- Actualidad} \\
\textsc{School of Health and Wellbeing -- University of Glasgow} & \textsc{Glasgow -- UK} \\
\multicolumn{2}{p{18cm}}{%
\begin{itemize}
  \item Estimé los efectos causales del desempleo sobre desenlaces de salud física y mental empleando métodos-g con datos administrativos.
  \item Realicé análisis de mediación causal del ingreso relativo entre el desempleo y la mortalidad por todas las causas usando datos administrativos enlazados.
  \item Modelé el efecto de implementar políticas económicas sobre desenlaces de salud mediante modelamiento por microsimulación.
\end{itemize}
} \\
\\
\textbf{Asistente de Investigación de Posgrado} & \textbf{Jul 2024 -- Jun 2025} \\
\textsc{Institute for Health Metrics and Evaluation} & \textsc{Seattle -- USA} \\
\multicolumn{2}{p{18cm}}{%
\begin{itemize}
  \item Desarrollé hiperparámetros para un modelo de meta-regresión jerárquico bayesiano con el fin de estimar prevalencia, incidencia y ocurrencia de desenlaces clínicos de infecciones del tracto respiratorio inferior en 80 países.
  \item Ejecuté cadenas MCMC utilizando priors basados en datos y propagación de la incertidumbre para estimar el traslape probabilístico entre grupos de inclusión social en 22 países de altos ingresos.
  \item Estimé fracciones atribuibles poblacionales y razones de mortalidad estandarizadas a partir del traslape entre grupos de inclusión social.
\end{itemize}
} \\
\\
\textbf{Asistente de Investigación de Posgrado} & \textbf{Jun 2024 -- Mar 2025} \\
\textsc{Department of Global Health -- University of Washington} & \textsc{Seattle -- USA} \\
\multicolumn{2}{p{18cm}}{%
\begin{itemize}
  \item Diseñé un protocolo de estudio para evaluar la efectividad de estrategias de implementación para la atención de quemaduras pediátricas en establecimientos de salud peruanos.
  \item Coordiné un estudio de métodos mixtos para explorar barreras y facilitadores en más de 10 hospitales peruanos de bajos recursos.
  \item Co-desarrollé estrategias de comunicación con 2 organizaciones de actores clave para difundir hallazgos complejos entre tomadores de decisión en salud.
\end{itemize}
} \\
\\
\textbf{Asistente de Investigación de Posgrado} & \textbf{Sep 2023 -- Mar 2024} \\
\textsc{Health Population Research Center -- University of Washington} & \textsc{Seattle -- USA} \\
\multicolumn{2}{p{18cm}}{%
\begin{itemize}
  \item Mantuve y almacené datos sociodemográficos y de utilización como parte de una encuesta de seguimiento posterior a un ensayo clínico en escolares, utilizando plataformas electrónicas de captura de datos.
  \item Traduje materiales de sometimiento al comité de ética (IRB) del español al inglés para el Alzheimer's Disease Research Center.
  \item Analicé datos de participación y alcance de ensayos clínicos para generar herramientas de visualización y tablas resumen, permitiendo que los gestores de programa ajustaran las estrategias de reclutamiento.
\end{itemize}
} \\
\\
\newpage
\textbf{Consultor de Investigación} & \textbf{Ene 2023 -- Ene 2024} \\
\textsc{Health Technology Assessment and Research Institute} & \textsc{Lima -- Perú} \\
\multicolumn{2}{p{18cm}}{%
\begin{itemize}
  \item Integré datos de vigilancia y datos clínicos de hospitales peruanos en modelos estadísticos de carga de enfermedad para apoyar los procesos nacionales de evaluación de tecnologías sanitarias y priorización.
  \item Realicé análisis de impacto presupuestario y modelos de costo-efectividad que priorizaron intervenciones oncológicas por montos superiores a USD 1 millón.
  \item Sinteticé los hallazgos en informes ejecutivos dirigidos a tomadores de decisión y socios, orientando la política de reembolso del Seguro Social de Salud del Perú.
\end{itemize}
} \\
\\
\textbf{Asistente de Investigación} & \textbf{Mar 2019 -- Ago 2023} \\
\textsc{CRONICAS Center of Excellence in Chronic Diseases} & \textsc{Lima -- Perú} \\
\multicolumn{2}{p{18cm}}{%
\begin{itemize}
  \item Coordiné 2 talleres por sede con agentes comunitarios de salud en 3 regiones peruanas para co-diseñar una aplicación móvil para la detección temprana del deterioro cognitivo.
  \item Contribuí a propuestas de financiamiento en epidemiología de enfermedades no transmisibles en poblaciones pediátricas y adultas mediante análisis preliminares de datos clínicos y demográficos, apoyando el desarrollo de proyectos de investigación.
  \item Fui coautor de 5 manuscritos científicos, incluyendo análisis de datos longitudinales, revisiones sistemáticas y modelos de aprendizaje automático.
\end{itemize}
} \\
\\
\end{longtable}
